\chapter{Roles, Permissões e Segurança em Banco de Dados}
\label{cap:roles_permissoes_seguranca}

Uma base de dados não é apenas um conjunto de estruturas técnicas. Ela também é um ambiente compartilhado, no qual diferentes pessoas e perfis interagem com os dados segundo necessidades distintas. Em um sistema real, nem todos os usuários devem ter o mesmo nível de acesso. Alguns precisam apenas consultar; outros precisam inserir ou atualizar dados; poucos devem possuir permissões amplas de administração.

É nesse contexto que a DCL (\textit{Data Control Language}) assume papel central. Esta aula aborda os mecanismos de controle de acesso em bancos de dados relacionais, com foco em roles, concessão de privilégios, revogação de permissões e boas práticas de segurança.

\section{Segurança como parte da implementação}

Em muitos projetos iniciantes, segurança é tratada como assunto secundário. Entretanto, do ponto de vista profissional, ela faz parte da própria implementação do sistema.

Em um ambiente acadêmico, por exemplo:
\begin{itemize}
    \item alunos podem visualizar apenas seus próprios dados;
    \item professores podem consultar turmas e lançar notas;
    \item secretarias podem cadastrar e atualizar matrículas;
    \item administradores podem criar estruturas e gerenciar acessos.
\end{itemize}

Conceder a todos o mesmo nível de acesso resultaria em risco de inconsistência, uso indevido e perda de controle sobre o banco. A DCL existe justamente para organizar essas fronteiras.

\section{O que é DCL}

A DCL, ou \textit{Data Control Language}, é o subconjunto da SQL voltado ao controle de permissões e acessos.

Seus principais comandos são:
\begin{itemize}
    \item \texttt{CREATE ROLE}
    \item \texttt{ALTER ROLE}
    \item \texttt{DROP ROLE}
    \item \texttt{GRANT}
    \item \texttt{REVOKE}
\end{itemize}

Esses comandos permitem definir quem pode acessar o banco, o que cada perfil pode fazer e sobre quais objetos esses privilégios recaem.

\section{Usuários, grupos e roles}

No PostgreSQL, a noção central é a de \texttt{role}. Uma role pode representar:
\begin{itemize}
    \item um usuário individual;
    \item um grupo de permissões;
    \item um perfil administrativo;
    \item uma entidade sem login, usada apenas para organização de privilégios.
\end{itemize}

Essa abordagem é útil porque permite separar a identidade da pessoa da lógica de permissões do sistema.

\section{Roles com e sem login}

\subsection{Role sem login}

Uma role sem login funciona como um perfil reutilizável.

\begin{sqlcode}
CREATE ROLE readonly;
CREATE ROLE secretaria_role;
CREATE ROLE admin_role;
\end{sqlcode}

Essas roles não representam necessariamente pessoas. Elas podem representar tipos de acesso.

\subsection{Role com login}

Uma role com login representa, na prática, um usuário autenticável.

\begin{sqlcode}
CREATE ROLE joao LOGIN PASSWORD 'senha_segura';
\end{sqlcode}

Nesse caso, \texttt{joao} pode acessar o banco usando autenticação.

\subsection{Role com parâmetros adicionais}

\begin{sqlcode}
CREATE ROLE maria
WITH
LOGIN
PASSWORD 'senha123'
VALID UNTIL '2026-12-31'
CONNECTION LIMIT 5;
\end{sqlcode}

Esse exemplo mostra que uma role pode ser criada com:
\begin{itemize}
    \item senha;
    \item validade temporal;
    \item limite de conexões simultâneas.
\end{itemize}

\section{Privilégios}

Privilégios são autorizações concedidas sobre objetos do banco de dados. Eles determinam o que uma role pode fazer.

Entre os privilégios mais comuns estão:
\begin{itemize}
    \item \texttt{SELECT}
    \item \texttt{INSERT}
    \item \texttt{UPDATE}
    \item \texttt{DELETE}
    \item \texttt{USAGE}
    \item \texttt{CREATE}
    \item \texttt{ALL PRIVILEGES}
\end{itemize}

Esses privilégios podem ser aplicados sobre:
\begin{itemize}
    \item tabelas;
    \item views;
    \item schemas;
    \item sequences;
    \item bancos de dados.
\end{itemize}

\section{GRANT}

O comando \texttt{GRANT} concede privilégios.

\subsection{Concedendo leitura}

\begin{sqlcode}
GRANT SELECT ON alunos TO readonly;
\end{sqlcode}

Agora a role \texttt{readonly} pode consultar a tabela \texttt{alunos}.

\subsection{Concedendo múltiplos privilégios}

\begin{sqlcode}
GRANT SELECT, INSERT, UPDATE ON matricula TO secretaria_role;
\end{sqlcode}

Nesse caso, a role \texttt{secretaria\_role} pode consultar, inserir e atualizar registros da tabela \texttt{matricula}.

\subsection{Concedendo permissão em schema}

\begin{sqlcode}
GRANT USAGE, CREATE ON SCHEMA academico TO admin_role;
\end{sqlcode}

Aqui:
\begin{itemize}
    \item \texttt{USAGE} permite utilizar o schema;
    \item \texttt{CREATE} permite criar objetos dentro dele.
\end{itemize}

\subsection{Concedendo role a usuário}

\begin{sqlcode}
GRANT readonly TO joao;
\end{sqlcode}

Esse comando associa ao usuário \texttt{joao} as permissões da role \texttt{readonly}.

\subsection{Concedendo privilégios amplos}

\begin{sqlcode}
GRANT ALL PRIVILEGES ON ALL TABLES IN SCHEMA public TO admin_role;
GRANT ALL PRIVILEGES ON DATABASE universidade TO admin_role;
\end{sqlcode}

Esses comandos são úteis em contextos administrativos, mas exigem cautela.

\section{REVOKE}

O comando \texttt{REVOKE} remove privilégios anteriormente concedidos.

\subsection{Revogando parte das permissões}

\begin{sqlcode}
REVOKE INSERT, UPDATE ON matricula FROM secretaria_role;
\end{sqlcode}

\subsection{Revogando tudo sobre uma tabela}

\begin{sqlcode}
REVOKE ALL PRIVILEGES ON alunos FROM readonly;
\end{sqlcode}

\subsection{Revogando em banco}

\begin{sqlcode}
REVOKE ALL PRIVILEGES ON DATABASE universidade FROM admin_role;
\end{sqlcode}

\section{ALTER ROLE}

O comando \texttt{ALTER ROLE} modifica características de uma role já existente.

\subsection{Alterando senha}

\begin{sqlcode}
ALTER ROLE maria WITH PASSWORD 'nova_senha';
\end{sqlcode}

\subsection{Alterando validade}

\begin{sqlcode}
ALTER ROLE maria VALID UNTIL '2027-12-31';
\end{sqlcode}

\subsection{Alterando limite de conexão}

\begin{sqlcode}
ALTER ROLE maria CONNECTION LIMIT 10;
\end{sqlcode}

\subsection{Concedendo SUPERUSER}

\begin{sqlcode}
ALTER ROLE maria WITH SUPERUSER;
\end{sqlcode}

\begin{danger}
O privilégio \texttt{SUPERUSER} deve ser reservado a administradores de confiança. Seu uso indiscriminado compromete a política de segurança do banco.
\end{danger}

\section{DROP ROLE}

O comando \texttt{DROP ROLE} remove uma role.

\begin{sqlcode}
DROP ROLE secretaria_role;
\end{sqlcode}

A remoção só será possível se não houver dependências incompatíveis, como privilégios ativos, objetos pertencentes à role ou vínculos de herança que impeçam a exclusão.

\section{Princípio do menor privilégio}

Uma das ideias mais importantes em segurança de banco de dados é o princípio do menor privilégio: cada usuário deve possuir apenas as permissões estritamente necessárias para exercer sua função.

Isso significa:
\begin{itemize}
    \item não conceder \texttt{DELETE} a quem só consulta;
    \item não conceder privilégios administrativos a quem apenas registra operações;
    \item não usar \texttt{SUPERUSER} quando privilégios limitados bastam;
    \item separar perfis operacionais de perfis técnicos.
\end{itemize}

\section{Consultando roles do sistema}

No PostgreSQL, a view \texttt{pg\_roles} pode ser usada para listar roles e algumas de suas características.

\begin{sqlcode}
SELECT * FROM pg_roles;
\end{sqlcode}

Esse recurso é útil para:
\begin{itemize}
    \item auditoria;
    \item conferência de perfis;
    \item entendimento do ambiente;
    \item verificação administrativa.
\end{itemize}

\section{Exemplo resolvido 1 -- Usuário somente leitura}

Suponha que um funcionário da coordenação precise apenas consultar dados acadêmicos.

\begin{sqlcode}
CREATE ROLE leitura_academica;
GRANT SELECT ON alunos TO leitura_academica;
GRANT SELECT ON cursos TO leitura_academica;
GRANT SELECT ON disciplinas TO leitura_academica;
\end{sqlcode}

Agora cria-se o usuário:

\begin{sqlcode}
CREATE ROLE paulo LOGIN PASSWORD 'senha123';
GRANT leitura_academica TO paulo;
\end{sqlcode}

\textbf{Interpretação:}
\begin{itemize}
    \item a role \texttt{leitura\_academica} concentra o perfil;
    \item o usuário \texttt{paulo} herda esse perfil;
    \item a manutenção fica mais simples do que conceder tudo diretamente ao usuário.
\end{itemize}

\section{Exemplo resolvido 2 -- Secretaria com manutenção restrita}

Suponha que a secretaria possa consultar e atualizar matrículas, mas não removê-las.

\begin{sqlcode}
CREATE ROLE secretaria_role;
GRANT SELECT, INSERT, UPDATE ON matricula TO secretaria_role;
REVOKE DELETE ON matricula FROM secretaria_role;
\end{sqlcode}

\textbf{Interpretação:}
\begin{itemize}
    \item a secretaria pode executar operações operacionais comuns;
    \item a remoção é bloqueada por ser mais sensível;
    \item isso protege o histórico contra exclusões indevidas.
\end{itemize}

\section{Exemplo resolvido 3 -- Grupo administrativo}

Considere um administrador responsável por objetos do schema acadêmico.

\begin{sqlcode}
CREATE ROLE admin_academico;
GRANT USAGE, CREATE ON SCHEMA academico TO admin_academico;
GRANT ALL PRIVILEGES ON ALL TABLES IN SCHEMA academico TO admin_academico;
\end{sqlcode}

Depois, cria-se um usuário e associa-se a role:

\begin{sqlcode}
CREATE ROLE renata LOGIN PASSWORD 'admin321';
GRANT admin_academico TO renata;
\end{sqlcode}

\textbf{Interpretação:}
\begin{itemize}
    \item a role concentra o perfil administrativo;
    \item a usuária herda o conjunto de privilégios;
    \item a administração permanece mais organizada.
\end{itemize}

\section{Exemplo resolvido 4 -- Revogando privilégio indevido}

Suponha que um operador tenha recebido permissão de atualização sobre a tabela \texttt{notas}, mas essa operação deva ser exclusiva dos professores.

\begin{sqlcode}
REVOKE UPDATE ON notas FROM operador_role;
\end{sqlcode}

\textbf{Interpretação:}
\begin{itemize}
    \item a política de acesso pode ser refinada ao longo do tempo;
    \item revogar é tão importante quanto conceder;
    \item segurança exige manutenção contínua.
\end{itemize}

\section{Exemplo resolvido 5 -- Role temporária}

Suponha a necessidade de criar um acesso temporário para auditoria.

\begin{sqlcode}
CREATE ROLE auditor_temporario
WITH LOGIN
PASSWORD 'auditoria'
VALID UNTIL '2025-12-31';
\end{sqlcode}

\begin{sqlcode}
GRANT SELECT ON alunos TO auditor_temporario;
GRANT SELECT ON matricula TO auditor_temporario;
\end{sqlcode}

\textbf{Interpretação:}
\begin{itemize}
    \item o acesso é limitado temporalmente;
    \item a role recebe apenas leitura;
    \item o risco operacional é menor do que em acessos amplos.
\end{itemize}

\section{Erro comum 1 -- Conceder privilégios amplos demais}

É um erro comum conceder \texttt{ALL PRIVILEGES} sem necessidade. Em muitos casos, o perfil requer apenas leitura ou atualização parcial.

\section{Erro comum 2 -- Confundir privilégio em banco com privilégio em tabela}

Ter acesso ao banco não significa automaticamente ter acesso a todas as tabelas. Os níveis de privilégio são diferentes e complementares.

\section{Erro comum 3 -- Usar SUPERUSER como solução fácil}

Transformar um usuário em \texttt{SUPERUSER} apenas para “fazer funcionar” é má prática e enfraquece completamente a política de segurança.

\section{Erro comum 4 -- Conceder privilégios diretamente a muitos usuários}

Quando muitos usuários recebem permissões individualmente, a manutenção fica difícil. O uso de roles de grupo é mais organizado e escalável.

\section{Erro comum 5 -- Não revisar permissões antigas}

Uma role criada para um contexto temporário pode continuar ativa sem necessidade. Revisões periódicas são essenciais.

\section{Boas práticas}

\begin{itemize}
    \item aplicar o princípio do menor privilégio;
    \item usar roles de grupo para perfis recorrentes;
    \item evitar privilégios administrativos sem necessidade;
    \item revisar permissões periodicamente;
    \item registrar a finalidade de cada role criada;
    \item preferir organização por perfis em vez de concessões diretas massivas.
\end{itemize}

\begin{quadro_dica}
Uma boa modelagem de segurança em banco de dados começa com a pergunta: “o que este usuário realmente precisa fazer?” e não “que privilégios posso conceder mais rapidamente?”.
\end{quadro_dica}

\section{Minicaso integrador -- Sistema acadêmico}

Considere um sistema acadêmico com quatro perfis:

\begin{itemize}
    \item \textbf{aluno}: consulta seus dados;
    \item \textbf{professor}: consulta turmas e lança notas;
    \item \textbf{secretaria}: gerencia matrículas;
    \item \textbf{administrador}: mantém a estrutura e as permissões.
\end{itemize}

Uma política inicial poderia ser:

\subsection*{Passo 1 -- Criar roles de perfil}

\begin{sqlcode}
CREATE ROLE role_aluno;
CREATE ROLE role_professor;
CREATE ROLE role_secretaria;
CREATE ROLE role_admin;
\end{sqlcode}

\subsection*{Passo 2 -- Conceder privilégios}

\begin{sqlcode}
GRANT SELECT ON alunos TO role_aluno;

GRANT SELECT ON turmas TO role_professor;
GRANT SELECT, UPDATE ON notas TO role_professor;

GRANT SELECT, INSERT, UPDATE ON matricula TO role_secretaria;

GRANT USAGE, CREATE ON SCHEMA academico TO role_admin;
GRANT ALL PRIVILEGES ON ALL TABLES IN SCHEMA academico TO role_admin;
\end{sqlcode}

\subsection*{Passo 3 -- Criar usuários}

\begin{sqlcode}
CREATE ROLE aluno01 LOGIN PASSWORD 'a1';
CREATE ROLE prof01 LOGIN PASSWORD 'p1';
CREATE ROLE sec01 LOGIN PASSWORD 's1';
CREATE ROLE admin01 LOGIN PASSWORD 'adm1';
\end{sqlcode}

\subsection*{Passo 4 -- Associar usuários aos perfis}

\begin{sqlcode}
GRANT role_aluno TO aluno01;
GRANT role_professor TO prof01;
GRANT role_secretaria TO sec01;
GRANT role_admin TO admin01;
\end{sqlcode}

\textbf{Comentário didático:}
Esse minicaso mostra que a segurança do banco pode ser tratada de forma estruturada e previsível. Em vez de conceder privilégios de forma dispersa, cria-se uma arquitetura de perfis, o que melhora manutenção e governança.

\section{Exercícios básicos}

\begin{enumerate}
    \item Crie uma role apenas para leitura.
    \item Crie um usuário com login e senha.
    \item Conceda permissão de \texttt{SELECT} sobre a tabela \texttt{alunos}.
    \item Revogue a permissão de \texttt{UPDATE} de uma role.
    \item Explique a diferença entre uma role com login e uma role sem login.
\end{enumerate}

\section{Exercícios intermediários}

\begin{enumerate}
    \item Crie uma role administrativa para um schema acadêmico.
    \item Associe um usuário a uma role de grupo.
    \item Escreva um exemplo com \texttt{GRANT} de múltiplos privilégios.
    \item Escreva um exemplo com \texttt{REVOKE ALL PRIVILEGES}.
    \item Explique por que privilégios em nível de banco não substituem permissões em nível de tabela.
\end{enumerate}

\section{Exercícios avançados}

\begin{enumerate}
    \item Modele um conjunto de roles para um sistema universitário com perfis de aluno, professor, secretaria e administrador.
    \item Escreva comandos SQL para conceder privilégios adequados a cada perfil.
    \item Explique quais riscos surgem quando operadores recebem privilégios administrativos.
    \item Dê um exemplo de cenário em que uma \texttt{VIEW} combinada com \texttt{GRANT SELECT} seja mais segura do que acesso direto à tabela.
\end{enumerate}

\section{Exercício integrador}

Considere um sistema acadêmico com as tabelas \texttt{alunos}, \texttt{cursos}, \texttt{notas} e \texttt{matricula}. Elabore uma política de acesso com:

\begin{enumerate}
    \item uma role de leitura;
    \item uma role operacional;
    \item uma role administrativa;
    \item um usuário vinculado à role operacional;
    \item comandos de concessão e revogação coerentes com esse cenário;
    \item uma justificativa para as decisões adotadas.
\end{enumerate}

\section{Síntese da aula}

Esta aula mostrou que a implementação de banco de dados inclui também controle de acesso, definição de perfis e gestão de privilégios. Foram estudados os comandos \texttt{CREATE ROLE}, \texttt{ALTER ROLE}, \texttt{DROP ROLE}, \texttt{GRANT} e \texttt{REVOKE}, além do raciocínio de segurança associado ao princípio do menor privilégio. Com exemplos resolvidos, erros comuns, exercícios progressivos e um minicaso integrador, o estudante passa a compreender que segurança em banco de dados é parte essencial da arquitetura do sistema. A próxima aula encerrará a Parte III com índices, transações e funções.